\chapter{Cadrage et Cahier des Charges}

Ce chapitre formalise les besoins fonctionnels et techniques du projet \textbf{My Smart Budget}. Il définit les objectifs mesurables, priorise les fonctionnalités selon la méthode MoSCoW et détaille l'architecture technique retenue pour répondre aux contraintes de performance et de sécurité.

\section{Objectifs métier, techniques et pédagogiques}

\subsection{Objectifs métier}

\textbf{My Smart Budget} répond à un besoin concret identifié auprès de plusieurs profils d'utilisateurs (étudiants, jeunes actifs, familles) : la difficulté à suivre les dépenses quotidiennes et à planifier un budget mensuel efficacement.

L'objectif principal est de \textbf{simplifier la gestion budgétaire} via une interface claire, automatisée et personnalisable, offrant une vision temps réel de la santé financière.

\subsubsection{Bénéfices attendus mesurables}

\begin{center}
\small
\begin{tabular}{|p{0.25\textwidth}|p{0.45\textwidth}|p{0.2\textwidth}|}
\hline
\textbf{Bénéfice} & \textbf{Indicateur de succès} & \textbf{Cible} \\
\hline
\textbf{Gain de temps} & Temps de gestion mensuelle & $< 15$ min \\
\hline
\textbf{Accessibilité} & Taux d'adoption (par persona) & $\geq 70\,\%$ \\
\hline
\textbf{Éducation} & Taux d'atteinte des objectifs d'épargne & $\geq 60\,\%$ \\
\hline
\textbf{Performance} & Temps d'accès au Dashboard & $< 5$ sec \\
\hline
\textbf{Fiabilité} & Réduction des erreurs vs Saisie manuelle & $-40\,\%$ \\
\hline
\end{tabular}
\end{center}

\subsection{Objectifs techniques}

Le projet repose sur une architecture \textbf{3-tiers évolutive} : Front-end (React 18), Back-end (Node.js 20/Express), Base de données (MongoDB 6).

\subsubsection{Exigences de performance (Non-fonctionnelles)}

\begin{center}
\small
\begin{tabular}{|p{0.15\textwidth}|p{0.3\textwidth}|p{0.15\textwidth}|p{0.3\textwidth}|}
\hline
\textbf{Critère} & \textbf{Métrique} & \textbf{Objectif} & \textbf{Justification} \\
\hline
\textbf{Performance} & Time-to-interactive & $< 2$s (P95) & UX Fluide \\
\hline
\textbf{Sécurité} & Auth (JWT + Bcrypt) & Cost $\geq 12$ & Protection RGPD \\
\hline
\textbf{Qualité} & Couverture de tests & $\geq 60\,\%$ & Maintenabilité \\
\hline
\textbf{Scalabilité} & Architecture Stateless & 10k users & Évolutivité \\
\hline
\textbf{Fiabilité} & Disponibilité (SLA) & $99.5\,\%$ & Accès continu \\
\hline
\end{tabular}
\end{center}

\subsection{Objectifs pédagogiques (Référentiel CDA)}

Ce projet permet de valider les trois blocs de compétences du Titre Professionnel de niveau 6.

\begin{itemize}
    \item \textbf{Concevoir et développer des composants d'interface :} Création d'un Design System, composants React réutilisables, respect normes WCAG 2.1.
    \item \textbf{Concevoir et développer la persistance des données :} Modélisation NoSQL (MongoDB), validation des schémas, scripts de migration.
    \item \textbf{Concevoir et développer une application multicouche :} Architecture API REST (OpenAPI), séparation des responsabilités, tests E2E.
\end{itemize}

\section{Priorisation des fonctionnalités (MoSCoW)}

Afin de garantir la livraison d'un MVP (Minimum Viable Product) fonctionnel dans les délais, les fonctionnalités ont été classées par priorité.

\subsection{Matrice de priorisation}

\begin{center}
\small
\begin{tabular}{|p{0.15\textwidth}|p{0.4\textwidth}|p{0.35\textwidth}|}
\hline
\textbf{Priorité} & \textbf{Fonctionnalité} & \textbf{Justification Métier} \\
\hline
\rowcolor{green!10} \textbf{Must Have} & Auth JWT \& Gestion de session & Sécurité et cloisonnement des données \\
\hline
\rowcolor{green!10} \textbf{Must Have} & CRUD Transactions \& Catégories & Cœur du système de suivi \\
\hline
\rowcolor{green!10} \textbf{Must Have} & Enveloppes Budgétaires & Contrôle des dépenses \\
\hline
\rowcolor{green!10} \textbf{Must Have} & Dashboard Synthétique & Aide à la décision immédiate \\
\hline
\textbf{Should Have} & Import CSV standardisé & Gain de temps critique \\
\hline
\textbf{Should Have} & Export Données (RGPD) & Conformité légale \\
\hline
\textbf{Could Have} & Notifications Seuils & Prévention des découverts \\
\hline
\textbf{Could Have} & Gestion Multi-comptes & Réalisme accru \\
\hline
\rowcolor{red!5} \textbf{Won't Have} & Agrégation Bancaire (DSP2) & Complexité technique et coût (v2) \\
\hline
\rowcolor{red!5} \textbf{Won't Have} & OCR Factures & Technologie IA/ML (v2) \\
\hline
\end{tabular}
\end{center}

\section{Périmètre du MVP (Version 1.0)}

\subsection{Fonctionnalités incluses}

\begin{description}
    \item[Authentification :] Inscription, connexion JWT (24h), hashage Bcrypt.
    \item[Gestion Profil :] Modification informations, changement mot de passe, suppression compte.
    \item[Transactions :] Saisie manuelle, catégorisation, filtres avancés.
    \item[Enveloppes :] Définition de plafonds par catégorie, calcul "reste à vivre".
    \item[Dashboard :] Graphiques dynamiques (Recharts), KPIs financiers.
    \item[Import/Export :] Traitement fichiers CSV, export JSON.
\end{description}

\begin{tcolorbox}[colback=yellow!5!white,colframe=yellow!60!black,title=\textbf{Suivi de projet GitHub — Données réelles}]
\textbf{État du Backlog (Janvier 2025) :}
\begin{itemize}
    \item \textbf{Total Issues :} 47 créées depuis le lancement.
    \item \textbf{Done :} 34 issues complétées (72\%).
    \item \textbf{In Progress :} 3 issues (Auth Refactor, Optim Dashboard).
    \item \textbf{Milestone MVP v1.0 :} 27/40 issues critiques terminées.
\end{itemize}

\textbf{Exemple d'Issue (US-02) :}
\textit{"En tant qu'utilisateur, je veux créer/éditer une transaction pour suivre mes dépenses."} \\
$\rightarrow$ Status : \textbf{Closed} | PR : \#127 | Effort : 8 pts.
\end{tcolorbox}

\section{Cibles et Personae}

Deux profils types ont guidé la conception de l'expérience utilisateur (UX).

\subsection{Persona 1 : Léa, 23 ans, Étudiante}
\textbf{Situation :} Revenus limités (800€), vit en colocation. \\
\textbf{Pain Points :} Fin de mois difficile, suivi Excel fastidieux. \\
\textbf{Besoin clé :} Visualiser son "reste à vivre" en moins de 5 secondes sur mobile avant un achat.

\subsection{Persona 2 : Marc, 38 ans, Père de famille}
\textbf{Situation :} Revenus confortables (4500€), gestion de budget familial complexe. \\
\textbf{Pain Points :} Dépassements sur le poste "Loisirs", manque de vision globale. \\
\textbf{Besoin clé :} Gérer des enveloppes budgétaires strictes et exporter des rapports mensuels.

\section{User Stories et Critères d'Acceptation}

\begin{tcolorbox}[colback=blue!5!white,colframe=blue!60!black,title=\textbf{Exemple de User Story critique : US-01 (Connexion)}]
\textbf{En tant que :} Utilisateur enregistré \\
\textbf{Je veux :} Me connecter à mon espace personnel \\
\textbf{Afin de :} Accéder à mes données financières de manière sécurisée

\textbf{Critères d'acceptation (Definition of Done) :}
\begin{enumerate}
    \item [$\boxtimes$] Le formulaire demande Email + Mot de passe.
    \item [$\boxtimes$] Le mot de passe est vérifié via \texttt{bcrypt.compare()}.
    \item [$\boxtimes$] Un Token JWT est généré (expiration 24h) au succès.
    \item [$\boxtimes$] Redirection vers \texttt{/dashboard}.
    \item [$\boxtimes$] Blocage après 3 tentatives échouées (Rate Limiting).
\end{enumerate}
\end{tcolorbox}

\section{Architecture Technique}

\subsection{Schéma fonctionnel 3-tiers}

\begin{center}
\small$\rightarrow$~\textit{Architecture 3-tiers — voir Annexe~\ref{ann:figures}, figure~\ref{fig:architecture-3-tiers}}
\end{center}

\subsection{Responsabilités des couches}

\subsubsection{Front-end (Client)}
\begin{itemize}
    \item \textbf{Tech :} React 18, Tailwind CSS, Recharts.
    \item \textbf{Rôle :} Interface utilisateur, validations formulaires (Yup), gestion de l'état (Zustand).
\end{itemize}

\subsubsection{Back-end (Serveur)}
\begin{itemize}
    \item \textbf{Tech :} Node.js, Express.js.
    \item \textbf{Rôle :} API REST, logique métier, authentification, sécurité.
\end{itemize}

\subsubsection{Data (Persistance)}
\begin{itemize}
    \item \textbf{Tech :} MongoDB, Mongoose ODM.
    \item \textbf{Rôle :} Stockage flexible, indexation (\texttt{userId + date}), agrégations statistiques.
\end{itemize}

\section{Choix Technologiques et Justifications}

\begin{center}
\small
\begin{tabular}{|p{0.2\textwidth}|p{0.2\textwidth}|p{0.2\textwidth}|p{0.3\textwidth}|}
\hline
\textbf{Domaine} & \textbf{Choix} & \textbf{Alternative} & \textbf{Raison du choix} \\
\hline
\textbf{Back-end} & \textbf{Express.js} & NestJS & Légèreté, flexibilité et courbe d'apprentissage rapide. \\
\hline
\textbf{Front-end} & \textbf{React.js} & Vue.js & Écosystème riche, composants réutilisables. \\
\hline
\textbf{Base de données} & \textbf{MongoDB} & PostgreSQL & Flexibilité du schéma JSON (NoSQL) adaptée au MVP. \\
\hline
\textbf{Validation} & \textbf{Joi / Yup} & Zod & Maturité et intégration native Express. \\
\hline
\textbf{Hébergement} & \textbf{Render} & AWS EC2 & Simplicité du déploiement (PaaS) et coût faible. \\
\hline
\end{tabular}
\end{center}

\begin{tcolorbox}[colback=green!5!white,colframe=green!60!black,title=\textbf{Validation technique — Métriques de production}]
Après 2 mois d'utilisation sur le périmètre MVP :
\begin{itemize}
    \item \textbf{MongoDB :} Temps moyen requête simple = 42ms.
    \item \textbf{API :} Temps d'agrégation dashboard = 156ms (pour 1800+ transactions).
    \item \textbf{Performance Front :} Score Lighthouse 92/100, Bundle size 890KB (Gzipped).
\end{itemize}
\end{tcolorbox}

\section{Validation Utilisateur et Bilan Qualité}

Une phase de test a été menée en novembre 2024 auprès de 5 utilisateurs représentatifs.

\subsection{Résultats quantitatifs}

\begin{center}
\begin{tabular}{|l|c|c|l|}
\hline
\textbf{Métrique} & \textbf{Résultat} & \textbf{Objectif} & \textbf{Statut} \\
\hline
Score SUS (Usability) & \textbf{78/100} & $\geq 75$ & \textcolor{green!60!black}{\textbf{Succès}} \\
\hline
Temps inscription & \textbf{1min 42s} & $< 2$min & \textcolor{green!60!black}{\textbf{Succès}} \\
\hline
Taux réussite (sans aide) & \textbf{84\,\%} & $\geq 80\,\%$ & \textcolor{green!60!black}{\textbf{Succès}} \\
\hline
Taux d'erreur & \textbf{7\,\%} & $< 10\,\%$ & \textcolor{green!60!black}{\textbf{Succès}} \\
\hline
\end{tabular}
\end{center}

\subsection{Améliorations suite aux retours}
\begin{enumerate}
    \item \textbf{Visibilité Import :} Le bouton d'import CSV a été déplacé dans le header (Feedback P1).
    \item \textbf{Terminologie :} Harmonisation des termes "Budget" et "Enveloppe".
    \item \textbf{Mobile :} Agrandissement des polices pour une meilleure lisibilité sur smartphone.
\end{enumerate}

\section{Roadmap et Perspectives}

Le projet suit une feuille de route claire pour les évolutions futures :

\begin{itemize}
    \item \textbf{v1.0 (Actuelle) :} MVP stable, déploiement Production.
    \item \textbf{v1.1 (Février 2025) :} Optimisation mobile, mode sombre, amélioration import CSV.
    \item \textbf{v2.0 (Avril 2025) :} POC d'agrégation bancaire automatique et notifications Push.
\end{itemize}

La validation de ces jalons confirme la capacité du projet à évoluer d'un prototype académique vers une solution exploitable par le grand public.